
% =====================================================================
\section{The real Dirac matrices of the \SIG{4+4}{7+1} spacetime}
\label{sec:matrices}
% =====================================================================

This section meets objective O1. It defines the eight matrices, proves that they are real and that
they satisfy the anticommutation relations of the metric, determines the group that their scaled
commutators generate, and derives the chiral projection operators. Throughout, $1_n$ denotes the
$n\times n$ unit matrix, a superscript $T$ the transpose, and a dagger the Hermitian conjugate. An
identifier such as \SIG{\texttt{44-C1}}{\texttt{71-C1}} names the machine check that verifies the
statement next to which it stands (Appendix~\ref{app:checks}).

%%<1644>
\subsection{Definition}
\label{sec:def}

The matrices are built in four steps, from $4\times4$ to $16\times16$ matrices. The $8\times8$ matrices
of the second and third steps are reduced generators of the kind that the author has used for the split
octonion algebra and for the spinors of $SO(4,4)$ \cite{Nash1990,Nash2010}; in those papers the index
takes the values $1,\dots,8$, and the value 8 there corresponds to the value 0 here.

\emph{Step 1: six antisymmetric $4\times4$ matrices.} Let $\epsilon_{hpq4}$ be the sign of the
permutation $(h,p,q,4)$ of $(1,2,3,4)$, and zero if two of the indices coincide. For $h=1,2,3$ and
$p,q=1,\dots,4$,
\begin{align}
(s_h)_{pq}&=\epsilon_{hpq4}-\big(\delta_{p4}\delta_{qh}-\delta_{ph}\delta_{q4}\big),\nonumber\\
(t_h)_{pq}&=\epsilon_{hpq4}+\big(\delta_{p4}\delta_{qh}-\delta_{ph}\delta_{q4}\big).
\label{eq:st}
\end{align}
The three $s_h$ are self-dual and the three $t_h$ anti-self-dual antisymmetric matrices; the letters
are those of Ref.~\cite{Nash1984}, where six matrices of this kind generate the rotations in a real
four-dimensional representation of the covering group of $SO(3,3)$.

\emph{Step 2: seven $8\times8$ matrices.} With $4\times4$ blocks,
\begin{align}
\tau^h&=\begin{pmatrix}0&s_h\\ s_h&0\end{pmatrix},\qquad
\tau^{7-h}=\begin{pmatrix}0&t_h\\ -t_h&0\end{pmatrix}\quad(h=1,2,3),\nonumber\\
\tau^7&=\tau^1\tau^2\tau^3\tau^4\tau^5\tau^6,\qquad\tau^0=1_8 .
\label{eq:tau}
\end{align}
One finds $\tau^7=\mathrm{diag}(-1_4,1_4)$ and
\begin{equation}
\tau^1\tau^2\tau^3=\tau^4\tau^5\tau^6\tau^7=\sigma_{88},\qquad
\sigma_{88}\equiv\begin{pmatrix}0&1_4\\ 1_4&0\end{pmatrix}
\label{eq:sigma88}
\end{equation}
(check \texttt{44-B1}).

\emph{Step 3: the barred matrices.} $\bar\tau^0=1_8$ and
$\bar\tau^A=\sigma_{88}\,(\tau^A)^{T}\sigma_{88}$ for $A=1,\dots,7$. The two sets satisfy
(check \texttt{44-B2})
\begin{equation}
\tau^A\bar\tau^B+\tau^B\bar\tau^A=2\,\eta^{AB}\,1_8=\bar\tau^A\tau^B+\bar\tau^B\tau^A,
\label{eq:taubar}
\end{equation}
where $\eta^{AB}$ is the metric defined in Eq.~(\ref{eq:clifford}) below.

\emph{Step 4: the eight $16\times16$ matrices.} In $8\times8$ blocks,
\begin{equation}
\gamma_{1644}^A=\begin{pmatrix}0&\bar\tau_{1644}^A\\ \tau_{1644}^A&0\end{pmatrix},\qquad A=0,1,\dots,7 .
\label{eq:gamma}
\end{equation}
The index $A$ is a contravariant index: these are the matrices that multiply $\partial_A$ in the Dirac
operator. Table~\ref{tab:gammas} of Appendix~\ref{app:matrices} displays all eight.

\subsection{Reality and the anticommutation relations}

\emph{Theorem 1.} (i) The eight matrices $\gamma_{1644}^A$ are real. Every entry is $0$, $+1$ or $-1$,
and every row and every column contains exactly one non-zero entry: each is a signed permutation
matrix. (ii) They satisfy
\begin{equation}
\gamma_{1644}^A\gamma_{1644}^B+\gamma_{1644}^B\gamma_{1644}^A=2\,\eta_{1644}^{AB}\,1_{16},
\label{eq:clifford}
\end{equation}
with the metric
\begin{equation}
\eta_{1644}=\mathrm{diag}(+1,+1,+1,+1,-1,-1,-1,-1).
\label{eq:eta}
\end{equation}

\emph{Proof.} (i) The matrices $s_h$ and $t_h$ are built from the integers $\epsilon_{hpq4}$ and
Kronecker deltas; $\tau^A$, $\bar\tau^A$ and $\gamma^A$ are built from them by forming blocks,
transposes and products. No complex number enters at any step (checks \texttt{44-R1},
\texttt{44-R2}). (ii) From the block form (\ref{eq:gamma}),
\begin{equation*}
\gamma^A\gamma^B+\gamma^B\gamma^A=\begin{pmatrix}\bar\tau^A\tau^B+\bar\tau^B\tau^A&0\\
0&\tau^A\bar\tau^B+\tau^B\bar\tau^A\end{pmatrix},
\end{equation*}
and both blocks equal $2\eta^{AB}1_8$ by Eq.~(\ref{eq:taubar}), an identity between integer matrices
that is verified for all 64 pairs of indices (checks \texttt{44-B2}, \texttt{44-C1}). $\blacksquare$

Thus $\gamma^0,\gamma^1,\gamma^2,\gamma^3$, which belong to the space-like coordinates, square to $+1$;
$\gamma^4,\gamma^5,\gamma^6,\gamma^7$, which belong to the time-like coordinates, square to $-1$; and
different matrices anticommute. Three consequences are used below (checks \texttt{44-C2} to
\texttt{44-C4}). First, $(\gamma^A)^T=\eta^{AA}\gamma^A$ (no sum): the space-like matrices are
symmetric and the time-like ones antisymmetric, and every $\gamma^A$ is an orthogonal matrix. Second,
the 256 ordered products $\gamma^{A_1}\cdots\gamma^{A_r}$ with $A_1<\dots<A_r$ and $r=0,\dots,8$ are
linearly independent. They therefore span all real $16\times16$ matrices: the Clifford algebra of the
metric (\ref{eq:eta}) is the full matrix algebra, $\Cl(4,4)\cong\mathbb{R}(16)$
\cite{AtiyahBottShapiro1964,LawsonMichelsohn1989}. Third, the 16-dimensional representation is
therefore irreducible over the real numbers.

\subsection{The scaled commutators and the group they generate}
\label{sec:pin}

The \emph{scaled commutators} are
\begin{equation}
\Sigma_{1644}^{AB}\equiv\tfrac14\big[\gamma_{1644}^A,\gamma_{1644}^B\big]
=\tfrac12\gamma_{1644}^A\gamma_{1644}^B\quad(A\neq B),
\label{eq:Sigma}
\end{equation}
where the second form follows from Eq.~(\ref{eq:clifford}). They are antisymmetric in their indices;
the 28 of them with $A<B$ are linearly independent (checks \texttt{44-L1}, \texttt{44-L2}).

\emph{Theorem 2.} (i) \emph{Lie algebra.} The scaled commutators satisfy
\begin{align}
[\Sigma^{AB},\Sigma^{CD}]={}&\eta^{BC}\Sigma^{AD}-\eta^{AC}\Sigma^{BD}\nonumber\\
&-\eta^{BD}\Sigma^{AC}+\eta^{AD}\Sigma^{BC},\label{eq:so44}\\
[\Sigma^{AB},\gamma^{C}]={}&\eta^{BC}\gamma^{A}-\eta^{AC}\gamma^{B}.\label{eq:vector}
\end{align}
Their real span is a Lie algebra isomorphic to $\so(4,4)$. (ii) \emph{Identity component.} The
matrices $S=\exp\big(\sum_{A<B}\theta_{AB}\Sigma^{AB}\big)$ with real parameters $\theta_{AB}$ generate
a connected group, $\Spin_0(4,4)$. Every such $S$ satisfies
\begin{equation}
S\gamma^AS^{-1}=\Lambda^A{}_B\gamma^B,\qquad\Lambda\,\eta\,\Lambda^T=\eta,\qquad\det\Lambda=+1,
\label{eq:lambda}
\end{equation}
with a real $8\times8$ matrix $\Lambda$. The map $S\mapsto\Lambda$ sends $\Spin_0(4,4)$ onto the
identity component of the pseudo-orthogonal group, and $S$ and $-S$ have the same image. (iii)
\emph{The whole group.} Each $\gamma^A$ is invertible, and the map
$\gamma^B\mapsto-\gamma^A\gamma^B(\gamma^A)^{-1}$ reverses $\gamma^A$ and leaves the other seven
matrices unchanged: it is the reflection of the axis $xA$. The group generated by $\Spin_0(4,4)$
together with the eight matrices is $\Pin(4,4)$. It has four connected components,
\begin{align}
\Pin(4,4)={}&\Spin_0(4,4)\,\cup\,\gamma^0\,\Spin_0(4,4)\nonumber\\
&\cup\,\gamma^4\,\Spin_0(4,4)\,\cup\,\gamma^0\gamma^4\,\Spin_0(4,4).
\label{eq:pin44}
\end{align}

In words: the scaled commutators span the Lie algebra of $\Pin(4,4)$ and generate, by exponentiation,
its identity component; the three other components are reached by multiplying with one space-like
matrix, with one time-like matrix, or with both. Exponentials alone can never leave the identity
component, because $\exp(sY)$, $0\le s\le1$, is a continuous path from the unit matrix to $\exp(Y)$.
This is why part (iii) is needed, and it is the precise sense in which the scaled commutators generate
$\Pin(4,4)$.

\emph{Proof.} (i) Equation (\ref{eq:vector}) follows from Eq.~(\ref{eq:clifford}):
$[\gamma^A\gamma^B,\gamma^C]=\gamma^A\{\gamma^B,\gamma^C\}-\{\gamma^A,\gamma^C\}\gamma^B
=2\eta^{BC}\gamma^A-2\eta^{AC}\gamma^B$. Applying it twice to $\Sigma^{CD}=\tfrac12\gamma^C\gamma^D$
gives Eq.~(\ref{eq:so44}); both are also verified for all index combinations (checks
\texttt{44-L3}, \texttt{44-L4}). Relations (\ref{eq:so44}) are the defining relations of the Lie
algebra of the pseudo-orthogonal group of the metric $\eta$. To identify the real form without appeal
to the literature, the program computes the Killing form
$B(X,Y)=\Tr(\mathrm{ad}_X\mathrm{ad}_Y)$ from the structure constants. It is diagonal in the basis
$\Sigma^{AB}$,
\begin{equation}
B(\Sigma^{AB},\Sigma^{AB})=-12\,\eta^{AA}\eta^{BB},
\label{eq:killing}
\end{equation}
and so has 12 negative and 16 positive eigenvalues (check \texttt{44-L5}). Equivalently, 12
generators are antisymmetric matrices, which exponentiate to rotations, and 16 are symmetric matrices,
which exponentiate to boosts (check \texttt{44-L6}). The numbers of compact and non-compact generators
of the real forms $\so(p,q)$ with $p+q=8$ are $(28,0)$ for $\so(8)$, $(21,7)$ for $\so(7,1)$,
$(16,12)$ for $\so(6,2)$, $(13,15)$ for $\so(5,3)$ and $(12,16)$ for $\so(4,4)$. The algebra is
$\so(4,4)$.

(ii) For a single plane the exponential is elementary, because
$(\gamma^A\gamma^B)^2=-\eta^{AA}\eta^{BB}$:
\begin{equation}
e^{\theta\Sigma^{AB}}=\begin{cases}\cos\frac\theta2+\gamma^A\gamma^B\sin\frac\theta2,&\eta^{AA}\eta^{BB}=+1,\\[3pt]
\cosh\frac\theta2+\gamma^A\gamma^B\sinh\frac\theta2,&\eta^{AA}\eta^{BB}=-1,\end{cases}
\label{eq:exp}
\end{equation}
a rotation in the first case and a boost in the second (check \texttt{44-G2}). By
Eq.~(\ref{eq:vector}) the derivative of $e^{sY}\gamma^Ce^{-sY}$ with respect to $s$ is again a linear
combination of the $\gamma$'s when $Y$ is a combination of the $\Sigma^{AB}$, so that
$S\gamma^CS^{-1}=\Lambda^C{}_D\gamma^D$; applying $S(\cdot)S^{-1}$ to Eq.~(\ref{eq:clifford}) shows
$\Lambda\eta\Lambda^T=\eta$. Check \texttt{44-G3} confirms both statements and $\det\Lambda=1$
numerically for random elements. A rotation by $\theta=2\pi$ in Eq.~(\ref{eq:exp}) gives $S=-1$ while
$\Lambda=1$: the map is two-to-one, and $-1$ is connected to $+1$ inside $\Spin_0(4,4)$.

(iii) $(\gamma^A)^{-1}=\eta^{AA}\gamma^A$ and the reflection property are check \texttt{44-G1}. By
definition $\Pin(4,4)$ is the group generated, inside the Clifford algebra, by the vectors
$v=v_A\gamma^A$ with $v^2=\pm1$ \cite{AtiyahBottShapiro1964,LawsonMichelsohn1989}. Every such $v$ is
obtained from $\gamma^0$ (if $v^2=+1$) or from $\gamma^4$ (if $v^2=-1$) by conjugation with an element
of $\Spin_0(4,4)$, so the group generated by $\Spin_0(4,4)$ and the $\gamma^A$ is $\Pin(4,4)$. A
product of two different matrices of the same kind is $\exp(\pi\Sigma^{AB})$ by Eq.~(\ref{eq:exp})
and lies in $\Spin_0(4,4)$; a product of a space-like and a time-like matrix reverses one space axis
and one time axis and does not. Hence the four components (\ref{eq:pin44}). $\blacksquare$

\subsection{Chirality: the real projection operators}
\label{sec:chirality}

Multiply the eight matrices in the order of their indices:
\begin{equation}
\omega_{1644}\equiv\gamma_{1644}^0\gamma_{1644}^1\cdots\gamma_{1644}^7=\begin{pmatrix}-1_8&0\\0&1_8\end{pmatrix}.
\label{eq:omega}
\end{equation}
Two properties of this \emph{volume element} follow from Eq.~(\ref{eq:clifford}) alone. First,
$\omega$ anticommutes with every $\gamma^A$: moving $\gamma^A$ through the product costs seven changes
of sign, one for each of the other matrices. Second,
$\omega^2=(-1)^{28}\prod_A(\gamma^A)^2=(+1)^4(-1)^4=+1$, where $(-1)^{28}$ is the sign of reversing the
order of eight anticommuting factors (checks \texttt{44-V1}, \texttt{44-V2}). Because $\omega$ is real
with $\omega^2=+1$, the matrices
\begin{align}
P_{1644}^{L}&=\tfrac12\big(1-\omega_{1644}\big)=\begin{pmatrix}1_8&0\\0&0\end{pmatrix},\nonumber\\
P_{1644}^{R}&=\tfrac12\big(1+\omega_{1644}\big)=\begin{pmatrix}0&0\\0&1_8\end{pmatrix}
\label{eq:PLPR}
\end{align}
are \emph{real} projectors. They satisfy $P^{L}+P^{R}=1$, $P^{L}P^{L}=P^{L}$, $P^{R}P^{R}=P^{R}$,
$P^{L}P^{R}=P^{R}P^{L}=0$ and $\Tr P^{L}=\Tr P^{R}=8$ (checks \texttt{44-P1}, \texttt{44-P2}). Two
further properties make them \emph{chiral} projectors. Every Dirac matrix turns a left-handed spinor
into a right-handed one and conversely, $\gamma^AP^{L}=P^{R}\gamma^A$ (check \texttt{44-P4}). And the
projectors commute with every $\Sigma^{AB}$ (check \texttt{44-P5}), so that the splitting
$\psi=P^{L}\psi+P^{R}\psi$ is the same in every Lorentz frame. In the present basis the splitting is
simply ``upper eight components, lower eight components.'' A real chirality is a property of the
$4+4$ signature; the companion paper shows that the $7+1$ signature has none.

\subsection{Other matrices that the theory needs}
\label{sec:other}

\emph{The two transposition matrices.} Let
\begin{align}
\Cm_{1644}&\equiv\gamma^0\gamma^1\gamma^2\gamma^3=\begin{pmatrix}-\sigma_{88}&0\\0&\sigma_{88}\end{pmatrix},\nonumber\\
\Cp_{1644}&\equiv\gamma^4\gamma^5\gamma^6\gamma^7=\begin{pmatrix}\sigma_{88}&0\\0&\sigma_{88}\end{pmatrix},
\label{eq:C}
\end{align}
so that $\omega=\Cm\Cp$ (checks \texttt{44-S1}, \texttt{44-S2}). Both are real, symmetric and square
to $1$, and
\begin{equation}
\Cm\gamma^A(\Cm)^{-1}=-(\gamma^A)^T,\qquad\Cp\gamma^A(\Cp)^{-1}=+(\gamma^A)^T
\label{eq:Ctranspose}
\end{equation}
(check \texttt{44-S3}). They play two roles: $\Cm$ defines the adjoint spinor of
Sec.~\ref{sec:dirac}, and both define the bilinear forms $\psi^TC^{\mp}\cdots\psi$ from which fermion
pairs are built (Sec.~\ref{sec:pairs}).

\emph{The complex structure.} The matrix $\gamma^4$ is real, squares to $-1$, and satisfies
$(\gamma^4)^T\Cm\gamma^4=\Cm$ (check \texttt{44-S5}). It will reappear as the generator of the
evolution of a particle at rest in the time $x4$.

\emph{The grading.} $W_{1644}\equiv i\,\Cm\gamma^4$ is Hermitian with $W^2=1$ and trace zero (checks
\texttt{44-D10}, \texttt{44-D11}). It appears in Sec.~\ref{sec:quantized}.

\emph{Relation to the matrices of the companion paper.} The eight real Dirac matrices
$\gamma_{1671}^A$ of the $7+1$ spacetime \cite{Companion1671} generate the same algebra of real
$16\times16$ matrices. They are related to the present ones by (check \texttt{71-B4})
\begin{align}
\gamma_{1671}^4&=\gamma_{1644}^0,\qquad\gamma_{1671}^j=-\gamma_{1644}^0\gamma_{1644}^j\quad(j=1,2,3),\nonumber\\
\gamma_{1671}^0&=+\gamma_{1644}^0\gamma_{1644}^5\gamma_{1644}^6\gamma_{1644}^7,\nonumber\\
\gamma_{1671}^5&=-\gamma_{1644}^0\gamma_{1644}^4\gamma_{1644}^6\gamma_{1644}^7,\nonumber\\
\gamma_{1671}^6&=+\gamma_{1644}^0\gamma_{1644}^4\gamma_{1644}^5\gamma_{1644}^7,\nonumber\\
\gamma_{1671}^7&=-\gamma_{1644}^0\gamma_{1644}^4\gamma_{1644}^5\gamma_{1644}^6 .
\label{eq:relation}
\end{align}
The matrix that the $4+4$ spacetime attaches to its hidden space-like axis is the time-like matrix of
the $7+1$ spacetime. The other seven matrices of $7+1$ are products of two or of four of the present
ones and are not linear combinations of them: the two sets are different sets of generators of the
same algebra.

\subsection{What has been established, and what has not}

Established, by exact arithmetic: the eight matrices are real; they satisfy the anticommutation
relations (\ref{eq:clifford}) of the $4+4$ metric; they generate the algebra of all real
$16\times16$ matrices; their scaled commutators span $\so(4,4)$, with 12 compact and 16 non-compact
generators, and generate, together with the matrices themselves, the four components of $\Pin(4,4)$;
the chirality matrix and the chiral projectors are real. The matrices rebuilt by the program are
identical, entry by entry, to those of the notebooks from which they were taken, and an independent
second program confirms the essential statements by other methods (Appendix~\ref{app:checks}).

Not established here, because it is not needed: the relation of these matrices to the triality of
$SO(4,4)$ and to the split octonions, for which we refer to Refs.~\cite{Nash1990,Nash2010}.
%%</1644>
%%<1671>
\subsection{Definition}
\label{sec:def}

The matrices are built in three steps. The first uses the real $8\times8$ matrices
$\tau_{1644}^A$ of the companion paper \cite{Companion1644}, which are reduced generators of the kind that the
author has used for the split octonion algebra \cite{Nash1990,Nash2010}; Appendix~\ref{app:matrices}
defines them, so that the present paper can be read by itself.

\emph{Step 1: seven octonionic $8\times8$ matrices.} Let
$\sigma_{88}$ be the $8\times8$ matrix with the $4\times4$ blocks $(0,1_4;1_4,0)$, and define, with
the labels of the present paper,
\begin{align}
\tau_{1671}^4&=1_8,\qquad\tau_{1671}^j=\tau_{1644}^j\quad(j=1,2,3),\nonumber\\
\tau_{1671}^0&=\tau_{1644}^4\,\sigma_{88},\qquad\tau_{1671}^k=\tau_{1644}^k\,\sigma_{88}\quad(k=5,6,7).
\label{eq:tau}
\end{align}
The seven matrices $\tau_{1671}^j$ with $j\in\{0,1,2,3,5,6,7\}$ are real and antisymmetric, they
anticommute with one another, and each squares to $-1_8$ (checks \texttt{71-B1}, \texttt{71-D6}):
\begin{equation}
\tau^i\tau^j+\tau^j\tau^i=-2\,\delta^{ij}\,1_8 .
\label{eq:tauclifford}
\end{equation}
They represent the multiplication by the seven imaginary units of the octonions
\cite{Bryant1982,HarveyLawson1982}.

\emph{Step 2: the time-like matrix.} In $8\times8$ blocks,
\begin{equation}
\gamma_{1671}^4=\begin{pmatrix}0&1_8\\1_8&0\end{pmatrix}.
\label{eq:gamma4}
\end{equation}

\emph{Step 3: the seven space-like matrices.}
\begin{equation}
\gamma_{1671}^j=\begin{pmatrix}-\tau_{1671}^j&0\\0&\tau_{1671}^j\end{pmatrix},\qquad j\in\{0,1,2,3,5,6,7\}
\label{eq:gamma}
\end{equation}
(check \texttt{71-B2}). The index $A$ of $\gamma_{1671}^A$ is a contravariant index: these are the
matrices that multiply $\partial_A$ in the Dirac operator. Table~\ref{tab:gammas} of
Appendix~\ref{app:matrices} displays all eight.

\emph{Relation to the matrices of the companion paper.} Both sets consist of real $16\times16$
matrices, and each generates the algebra of all such matrices. They are related by
(check \texttt{71-B4})
\begin{align}
\gamma_{1671}^4&=\gamma_{1644}^0,\qquad\gamma_{1671}^j=-\gamma_{1644}^0\gamma_{1644}^j\quad(j=1,2,3),\nonumber\\
\gamma_{1671}^0&=+\gamma_{1644}^0\gamma_{1644}^5\gamma_{1644}^6\gamma_{1644}^7,\nonumber\\
\gamma_{1671}^5&=-\gamma_{1644}^0\gamma_{1644}^4\gamma_{1644}^6\gamma_{1644}^7,\nonumber\\
\gamma_{1671}^6&=+\gamma_{1644}^0\gamma_{1644}^4\gamma_{1644}^5\gamma_{1644}^7,\nonumber\\
\gamma_{1671}^7&=-\gamma_{1644}^0\gamma_{1644}^4\gamma_{1644}^5\gamma_{1644}^6 .
\label{eq:relation}
\end{align}
The time-like matrix of the $7+1$ spacetime is the matrix that the $4+4$ spacetime attaches to its
hidden space-like axis. The other seven are products of two or of four of the $4+4$ matrices and are
not linear combinations of them: the two sets are different sets of generators of the same algebra.

\subsection{Reality and the anticommutation relations}

\emph{Theorem 1.} (i) The eight matrices $\gamma_{1671}^A$ are real. Every entry is $0$, $+1$ or $-1$,
and every row and every column contains exactly one non-zero entry: each is a signed permutation
matrix. (ii) They satisfy
\begin{equation}
\gamma_{1671}^A\gamma_{1671}^B+\gamma_{1671}^B\gamma_{1671}^A=2\,\eta_{1671}^{AB}\,1_{16},
\label{eq:clifford}
\end{equation}
with the metric
\begin{equation}
\eta_{1671}=\mathrm{diag}(-1,-1,-1,-1,+1,-1,-1,-1).
\label{eq:eta}
\end{equation}

\emph{Proof.} (i) The matrices are products and blocks of integer matrices; no complex number enters
(checks \texttt{71-R1}, \texttt{71-R2}). (ii) From the block forms (\ref{eq:gamma4}) and
(\ref{eq:gamma}): $(\gamma^4)^2=1$; $(\gamma^j)^2=\mathrm{diag}((\tau^j)^2,(\tau^j)^2)=-1$; for
$i\neq j$ the anticommutator of $\gamma^i$ and $\gamma^j$ is
$\mathrm{diag}(\tau^i\tau^j+\tau^j\tau^i,\ \tau^i\tau^j+\tau^j\tau^i)=0$ by
Eq.~(\ref{eq:tauclifford}); and
\begin{equation*}
\gamma^4\gamma^j+\gamma^j\gamma^4=\begin{pmatrix}0&\tau^j\\-\tau^j&0\end{pmatrix}+\begin{pmatrix}0&-\tau^j\\ \tau^j&0\end{pmatrix}=0 .
\end{equation*}
Equation (\ref{eq:clifford}) is also verified directly for all 64 pairs of indices (check
\texttt{71-C1}). $\blacksquare$

Thus the time-like matrix $\gamma^4$ squares to $+1$ and the seven space-like matrices square to
$-1$.

\emph{Why this sign convention is forced.} In the companion paper the space-like matrices square to
$+1$. Here that convention is impossible. Eight real $16\times16$ matrices with seven squares $+1$
and one square $-1$ do not exist, because the Clifford algebra with those signs is
$\Cl(7,1)\cong\mathbb{H}(8)$, the algebra of $8\times8$ matrices over the quaternions, whose smallest
real representation has 32 dimensions; whereas with one square $+1$ and seven squares $-1$ the algebra
is $\Cl(1,7)\cong\mathbb{R}(16)$ \cite{AtiyahBottShapiro1964,LawsonMichelsohn1989}. For the signature
$4+4$ the two conventions are equivalent. The consequence for the Dirac equation is the factor $i$ in
Eq.~(\ref{eq:dirac}) below, which is absent in the $4+4$ case.

Three consequences of Theorem 1 are used below (checks \texttt{71-C2} to \texttt{71-C4}). First,
$(\gamma^A)^T=\eta^{AA}\gamma^A$ (no sum): $\gamma^4$ is symmetric and the seven $\gamma^j$ are
antisymmetric, and every $\gamma^A$ is an orthogonal matrix. Second, the 256 ordered products
$\gamma^{A_1}\cdots\gamma^{A_r}$ with $A_1<\dots<A_r$ and $r=0,\dots,8$ are linearly independent, so
that the matrices generate all real $16\times16$ matrices. Third, the 16-dimensional representation is
therefore irreducible over the real numbers.

\subsection{The scaled commutators and the group they generate}
\label{sec:pin}

The \emph{scaled commutators} are
\begin{equation}
\Sigma_{1671}^{AB}\equiv\tfrac14\big[\gamma_{1671}^A,\gamma_{1671}^B\big]
=\tfrac12\gamma_{1671}^A\gamma_{1671}^B\quad(A\neq B),
\label{eq:Sigma}
\end{equation}
where the second form follows from Eq.~(\ref{eq:clifford}). They are antisymmetric in their indices;
the 28 of them with $A<B$ are linearly independent (checks \texttt{71-L1}, \texttt{71-L2}).

\emph{Theorem 2.} (i) \emph{Lie algebra.} The scaled commutators satisfy
\begin{align}
[\Sigma^{AB},\Sigma^{CD}]={}&\eta^{BC}\Sigma^{AD}-\eta^{AC}\Sigma^{BD}\nonumber\\
&-\eta^{BD}\Sigma^{AC}+\eta^{AD}\Sigma^{BC},\label{eq:so44}\\
[\Sigma^{AB},\gamma^{C}]={}&\eta^{BC}\gamma^{A}-\eta^{AC}\gamma^{B}.\label{eq:vector}
\end{align}
Their real span is a Lie algebra isomorphic to $\so(7,1)$. (ii) \emph{Identity component.} The
matrices $S=\exp\big(\sum_{A<B}\theta_{AB}\Sigma^{AB}\big)$ with real parameters $\theta_{AB}$ generate
a connected group, $\Spin_0(7,1)$. Every such $S$ satisfies
\begin{equation}
S\gamma^AS^{-1}=\Lambda^A{}_B\gamma^B,\qquad\Lambda\,\eta\,\Lambda^T=\eta,\qquad\det\Lambda=+1,
\label{eq:lambda}
\end{equation}
with a real $8\times8$ matrix $\Lambda$. The map $S\mapsto\Lambda$ sends $\Spin_0(7,1)$ onto the
proper orthochronous Lorentz group $SO_0(7,1)$, and $S$ and $-S$ have the same image. (iii)
\emph{The whole group.} Each $\gamma^A$ is invertible, and the map
$\gamma^B\mapsto-\gamma^A\gamma^B(\gamma^A)^{-1}$ reverses $\gamma^A$ and leaves the other seven
matrices unchanged: it is the reflection of the axis $xA$. The group generated by $\Spin_0(7,1)$
together with the eight matrices is the pin group of the $7+1$ spacetime. It has four connected
components,
\begin{align}
\Pin(7,1)={}&\Spin_0(7,1)\,\cup\,\gamma^1\,\Spin_0(7,1)\nonumber\\
&\cup\,\gamma^4\,\Spin_0(7,1)\,\cup\,\gamma^1\gamma^4\,\Spin_0(7,1).
\label{eq:pin44}
\end{align}

In words: the scaled commutators span the Lie algebra of $\Pin(7,1)$ and generate, by exponentiation,
its identity component; the three other components are reached by multiplying with one space-like
matrix, with the time-like matrix, or with both. Exponentials alone can never leave the identity
component, because $\exp(sY)$, $0\le s\le1$, is a continuous path from the unit matrix to $\exp(Y)$.
This is why part (iii) is needed, and it is the precise sense in which the scaled commutators generate
$\Pin(7,1)$.

\emph{Proof.} (i) Equation (\ref{eq:vector}) follows from Eq.~(\ref{eq:clifford}):
$[\gamma^A\gamma^B,\gamma^C]=\gamma^A\{\gamma^B,\gamma^C\}-\{\gamma^A,\gamma^C\}\gamma^B
=2\eta^{BC}\gamma^A-2\eta^{AC}\gamma^B$. Applying it twice to $\Sigma^{CD}=\tfrac12\gamma^C\gamma^D$
gives Eq.~(\ref{eq:so44}); both are also verified for all index combinations (checks
\texttt{71-L3}, \texttt{71-L4}). Relations (\ref{eq:so44}) are the defining relations of the Lie
algebra of the pseudo-orthogonal group of the metric $\eta$. To identify the real form without appeal
to the literature, the program computes the Killing form
$B(X,Y)=\Tr(\mathrm{ad}_X\mathrm{ad}_Y)$ from the structure constants. It is diagonal in the basis
$\Sigma^{AB}$,
\begin{equation}
B(\Sigma^{AB},\Sigma^{AB})=-12\,\eta^{AA}\eta^{BB},
\label{eq:killing}
\end{equation}
and so has 21 negative and 7 positive eigenvalues (check \texttt{71-L5}). Equivalently, 21 generators
are antisymmetric matrices, the rotations among the seven space coordinates, and 7 are symmetric
matrices, the boosts $\Sigma^{4j}$ (check \texttt{71-L6}). The numbers of compact and non-compact
generators of the real forms $\so(p,q)$ with $p+q=8$ are $(28,0)$ for $\so(8)$, $(21,7)$ for
$\so(7,1)$, $(16,12)$ for $\so(6,2)$, $(13,15)$ for $\so(5,3)$ and $(12,16)$ for $\so(4,4)$. The
algebra is $\so(7,1)$.

(ii) For a single plane the exponential is elementary, because
$(\gamma^A\gamma^B)^2=-\eta^{AA}\eta^{BB}$:
\begin{equation}
e^{\theta\Sigma^{AB}}=\begin{cases}\cos\frac\theta2+\gamma^A\gamma^B\sin\frac\theta2,&\eta^{AA}\eta^{BB}=+1,\\[3pt]
\cosh\frac\theta2+\gamma^A\gamma^B\sinh\frac\theta2,&\eta^{AA}\eta^{BB}=-1,\end{cases}
\label{eq:exp}
\end{equation}
a rotation in the first case and a boost in the second (check \texttt{71-G2}). By
Eq.~(\ref{eq:vector}) the derivative of $e^{sY}\gamma^Ce^{-sY}$ with respect to $s$ is again a linear
combination of the $\gamma$'s when $Y$ is a combination of the $\Sigma^{AB}$, so that
$S\gamma^CS^{-1}=\Lambda^C{}_D\gamma^D$; applying $S(\cdot)S^{-1}$ to Eq.~(\ref{eq:clifford}) shows
$\Lambda\eta\Lambda^T=\eta$. Check \texttt{71-G3} confirms both statements and $\det\Lambda=1$
numerically for random elements. A rotation by $\theta=2\pi$ in Eq.~(\ref{eq:exp}) gives $S=-1$ while
$\Lambda=1$: the map is two-to-one, and $-1$ is connected to $+1$ inside $\Spin_0(7,1)$.

(iii) $(\gamma^A)^{-1}=\eta^{AA}\gamma^A$ and the reflection property are check \texttt{71-G1}. By
definition the pin group is the group generated, inside the Clifford algebra, by the vectors
$v=v_A\gamma^A$ with $v^2=\pm1$ \cite{AtiyahBottShapiro1964,LawsonMichelsohn1989}. Every such $v$ is
obtained from $\gamma^4$ (if $v^2=+1$) or from $\gamma^1$ (if $v^2=-1$) by conjugation with an element
of $\Spin_0(7,1)$, so the group generated by $\Spin_0(7,1)$ and the $\gamma^A$ is the pin group. A
product of two different space-like matrices is $\exp(\pi\Sigma^{ij})$ by Eq.~(\ref{eq:exp}) and lies
in $\Spin_0(7,1)$; a product of a space-like matrix and the time-like one reverses one space axis and
the time axis and does not. Hence the four components (\ref{eq:pin44}). $\blacksquare$

\emph{What the name $\Pin(7,1)$ means.} We use it for the pin group of the spacetime with seven
space-like dimensions and one time-like dimension, realized by the real matrices $\gamma_{1671}^A$.
Two remarks make this precise. First, the matrices generate the Clifford algebra $\Cl(1,7)$, in which
the time-like generator squares to $+1$; in the notation that records these signs the group of
Theorem 2 is written $\Pin(1,7)$. The full Lorentz group $O(7,1)$ has a second double cover, the pin
group of the algebra $\Cl(7,1)$ with the opposite signs, in which a reflection of a space axis squares
to $+1$ instead of $-1$. Its generators would be eight real $16\times16$ matrices with seven squares
$+1$ and one square $-1$, and such matrices do not exist. The two groups have the same identity
component and the same Lie algebra and differ only in the components that contain reflections. With
purely real $16\times16$ Dirac matrices the group of Theorem 2 is the only possibility. Second, it is
not the group of the companion paper: $\so(7,1)$ has 21 compact and 7 non-compact generators,
$\so(4,4)$ has 12 and 16, and the signature of the Killing form (\ref{eq:killing}) is an invariant
that distinguishes them, although both sets of matrices live in the same algebra of real
$16\times16$ matrices.

\subsection{Chirality: no real projector, and the complex projectors}
\label{sec:chirality}

Multiply the eight matrices in the order of their indices:
\begin{equation}
\omega_{1671}\equiv\gamma_{1671}^0\gamma_{1671}^1\cdots\gamma_{1671}^7=\begin{pmatrix}0&-1_8\\1_8&0\end{pmatrix}.
\label{eq:omega}
\end{equation}
This \emph{volume element} anticommutes with every $\gamma^A$, because moving $\gamma^A$ through the
product costs seven changes of sign; but now
$\omega^2=(-1)^{28}\,(+1)(-1)^7=-1$ (checks \texttt{71-V1}, \texttt{71-V2}). A chirality matrix must
square to $+1$.

\emph{Theorem 3 (no real chirality in $7+1$).} A real $16\times16$ matrix that anticommutes with all
eight $\gamma_{1671}^A$ is a real multiple of $\omega_{1671}$. Since $\omega_{1671}^2=-1$, no real
matrix with square $+1$ anticommutes with all of them, and no real chiral projector exists.

\emph{Proof.} Expand a real matrix $X$ in the 256 products $\gamma^{A_1}\cdots\gamma^{A_r}$. A product
of rank $r$ anticommutes with $\gamma^A$ if $A$ is among its indices and $r$ is even, or if $A$ is not
among them and $r$ is odd; it commutes otherwise. A product therefore anticommutes with all eight
matrices only if every index occurs and $r$ is even, that is, only for $r=8$. Moreover the map
$X\mapsto X\gamma^A+\gamma^AX$ sends each product either to zero or to $\pm2$ times another product,
and different products to different products, so that the terms of a general $X$ cannot cancel one
another. Hence $X$ anticommutes with all eight matrices only if it is a multiple of the product of
rank 8, which is $\omega$, and the square of $c\,\omega$ is $-c^2\le0$. The program confirms the
statement in a second way: the linear system $X\gamma^A+\gamma^AX=0$ for the 256 entries of $X$ has
rank 255 (check \texttt{71-V3}). $\blacksquare$

With complex numbers a chirality matrix exists,
\begin{equation}
\chi_{1671}\equiv i\,\omega_{1671}=\begin{pmatrix}0&-i\,1_8\\ i\,1_8&0\end{pmatrix},\quad\chi^\dagger=\chi,\quad\chi^2=1,
\label{eq:chi}
\end{equation}
which anticommutes with every $\gamma^A$ (check \texttt{71-P1}). The chiral projectors are
\begin{align}
P_{1671}^{L}&=\tfrac12\big(1-\chi_{1671}\big)=\tfrac12\begin{pmatrix}1_8&i\,1_8\\-i\,1_8&1_8\end{pmatrix},\nonumber\\
P_{1671}^{R}&=\tfrac12\big(1+\chi_{1671}\big)=\tfrac12\begin{pmatrix}1_8&-i\,1_8\\ i\,1_8&1_8\end{pmatrix}.
\label{eq:PLPR}
\end{align}
They are Hermitian and satisfy $P^{L}+P^{R}=1$, $P^{L}P^{L}=P^{L}$, $P^{R}P^{R}=P^{R}$,
$P^{L}P^{R}=0$ and $\Tr P^{L}=\Tr P^{R}=8$; every Dirac matrix turns a left-handed spinor into a
right-handed one, $\gamma^AP^{L}=P^{R}\gamma^A$; and they commute with every $\Sigma^{AB}$, so that the
splitting $\psi=P^{L}\psi+P^{R}\psi$ is the same in every Lorentz frame (checks \texttt{71-P2} to
\texttt{71-P4}). The Dirac matrices remain real; only the projectors are complex.

The real projectors of the $4+4$ spacetime, $P_{1644}^{L}=\mathrm{diag}(1_8,0)$ and
$P_{1644}^{R}=\mathrm{diag}(0,1_8)$, are not chiral projectors for the present matrices, by
Theorem 3. They are, however, unitarily equivalent to the complex ones: with the unitary matrix
$U_{1671}^{c}=(1-i\,\gamma_{1671}^4)/\sqrt2$,
\begin{equation}
U^{c}P_{1671}^{L}(U^{c})^{-1}=P_{1644}^{R},\qquad U^{c}P_{1671}^{R}(U^{c})^{-1}=P_{1644}^{L}
\label{eq:Uc}
\end{equation}
(check \texttt{71-P5}); which of the two is called left-handed is a convention.

\subsection{Other matrices that the theory needs}
\label{sec:other}

\emph{The adjoint matrix} is the time-like matrix itself: $\gamma_{1671}^4$ is real, symmetric, squares
to $1$, and satisfies $(\gamma^A)^T=\gamma^4\gamma^A\gamma^4$ (check \texttt{71-S1}). The adjoint
spinor is $\bar\psi\equiv\psi^\dagger\gamma^4$.

\emph{The two transposition matrices} (check \texttt{71-K0}),
\begin{equation}
\Cp_{1671}\equiv\gamma_{1671}^4,\qquad\Cm_{1671}\equiv\gamma_{1671}^4\,\omega_{1671}=\begin{pmatrix}1_8&0\\0&-1_8\end{pmatrix},
\label{eq:C}
\end{equation}
are real and symmetric and satisfy
\begin{equation}
C^{\pm}\gamma^A(C^{\pm})^{-1}=\pm(\gamma^A)^T .
\label{eq:Ctranspose}
\end{equation}
They define the bilinear forms $\psi^TC^{\mp}\cdots\psi$ from which fermion pairs are built
(Sec.~\ref{sec:pairs}).

\emph{The Hamiltonian matrices.} $\beta_{1671}\equiv\gamma^4$ and $\alpha_{1671}^j\equiv\gamma^4\gamma^j$
for the seven space indices $j$ are real and symmetric, with
$\alpha^i\alpha^j+\alpha^j\alpha^i=2\delta^{ij}$, $\alpha^j\beta+\beta\alpha^j=0$ and $\beta^2=1$
(check \texttt{71-D4}). In blocks,
\begin{equation}
\beta_{1671}=\begin{pmatrix}0&1_8\\1_8&0\end{pmatrix},\qquad\alpha_{1671}^j=\begin{pmatrix}0&\tau_{1671}^j\\-\tau_{1671}^j&0\end{pmatrix}.
\label{eq:alphabeta}
\end{equation}

\subsection{What has been established, and what has not}

Established, by exact arithmetic: the eight matrices are real; they satisfy the anticommutation
relations (\ref{eq:clifford}) of the $7+1$ metric, in the only sign convention that real matrices
allow; they generate the algebra of all real $16\times16$ matrices; their scaled commutators span
$\so(7,1)$, with 21 compact and 7 non-compact generators, and generate, together with the matrices
themselves, the four components of $\Pin(7,1)$; no real chirality matrix exists, and the chiral
projectors are the complex matrices (\ref{eq:PLPR}). The matrices rebuilt by the program are
identical, entry by entry, to those of the notebooks from which they were taken, and an independent
second program confirms the essential statements by other methods (Appendix~\ref{app:checks}).

Not established here, because it is not needed: the relation of these matrices to the octonions and
to the exceptional geometries of seven and eight dimensions, for which we refer to
Refs.~\cite{Bryant1982,HarveyLawson1982,BrownGray1967}.
%%</1671>
